\documentclass[10pt,a4paper]{amsart}
\usepackage[margin=19mm]{geometry}
\usepackage{amsmath,amssymb,mathtools}
\usepackage[scr=rsfs,cal=euler]{mathalfa}
\usepackage{xfrac}
\usepackage[unicode,hidelinks,bookmarks=false]{hyperref}
\newcommand{\ie}{i.e.\ }
\newcommand{\cf}{cf.\ }
\newcommand{\resp}{respectively\ }
\newcommand{\Subsection}{\subsection}
\providecommand{\nobreakdash}{\nobreak\hbox{-}}
\setlength{\emergencystretch}{2em}
\multlinegap0pt
\usepackage{bm}
\usepackage{bbm}
\usepackage{dsfont}
\usepackage{xspace}
\usepackage{cancel}
\usepackage{mathtools}
\usepackage{amsthm}
\makeatletter
\@ifundefined{proposition}{\newtheorem{proposition}{Proposition}}{}
\makeatother
\newcommand{\ninv}{\mathord{\sim}}  
\title[Degree-preserving Godel logics with an involution: intermedi]{Degree-preserving Godel logics with an involution: intermediate logics and (ideal) paraconsistency}
\author{M.E. Coniglio, F. Esteva, J. Gispert, L. Godo}
\date{Source: arXiv:2601.08474v1 (CC BY 4.0)}
\begin{document}
\maketitle

\noindent\emph{Source and licence.} Degree-preserving Godel logics with an involution: intermediate logics and (ideal) paraconsistency. Version arXiv:2601.08474v1 (CC BY 4.0), distributed under the Creative Commons Attribution 4.0 International licence, as recorded in the arXiv licence metadata for this exact version and shown on the abstract page https://arxiv.org/abs/2601.08474v1. Full rights notice: licenses/arxiv-2601.08474-CC-BY-4.0.txt.\par\smallskip

\noindent\textit{Editorial scope.} Counted unit: the Proposition whose source label is \texttt{None} in the article (source0.tex lines 424--447, source label \texttt{None}), delivered together with the article's own complete original proof of it (source0.tex lines 449--511, one \texttt{proof} environment of 63 lines). No mathematics is written, repaired or re-derived: statement and proof are the article's own text line for line. Required context retained uncounted: none. Every definition, display and label of those lines is retained unchanged. Omitted from this package and not redistributed: 1--423, 448--448, 512--1188. Nothing inside a retained excerpt is omitted or altered: every retained body is delivered exactly as the article's lines are written, so no omission receipt of any kind accompanies this package. Citations: the retained excerpts carry no citation command at all, so no bibliography entry of the article is transcribed and no reference is redistributed. Reference adaptation: none was needed and none was made; every reference inside the delivered unit resolves inside it, so no unresolved reference or citation is produced. Printed slips kept literal: no printed word, symbol, hypothesis or number of the counted statement or of its proof was altered. Layout and class: the article's own class is \texttt{article} with packages including latexsym, amssymb, amsmath, amsthm, amsfonts, graphicx, color; the wrapper is the lane's pinned \texttt{amsart} editorial wrapper at 10pt on A4, which restates the theorem-like environments the retained text needs and adds the article's own macro definitions that the retained text needs (\texttt{\textbackslash ninv}), with extra wrapper packages (bm, bbm, dsfont, xspace, cancel, mathtools). The delivered numbering is the unit's own. Assets: no retained span contains a figure environment, a graphics-inclusion command, a PDF-page inclusion, a table or a TikZ picture; no image file is copied into this package, the package declares no asset dependency and no image file is redistributed with it. Licence: version arXiv:2601.08474v1 of this article is distributed under the Creative Commons Attribution 4.0 International licence, as recorded in the arXiv licence metadata for this exact version and shown on the abstract page https://arxiv.org/abs/2601.08474v1. A full rights notice is delivered at licenses/arxiv-2601.08474-CC-BY-4.0.txt. Characters: every retained excerpt is pure ASCII, so the delivered engine, XeLaTeX with its default Latin Modern text font, reports no missing character.
\begin{proposition}
The logics ${\rm G}^{[a}_\sim = \langle [0,1]_{\rm{G}_\sim}, F_{[a}\rangle$ for $a \in (0, 1]$,  ${\rm G}^{(a}_\sim=\langle [0,1]_{\rm{G}_\sim}, F_{(a} \rangle$ for $a \in [0, 1)$, and their finitary companions, satisfy the following properties:\footnote{In the following we use $p$ and $n$ to denote positive and negative values in $[0, 1]$ with respect to the negation $\ninv x = 1-x$; in other words, $p > 1/2$ and $n < 1/2$.}
\begin{enumerate}
\item[P1.] $\vdash_{[p} \; = \; \vdash_{[p'}$ and $\vdash_{(p} \; = \; \vdash_{(p'}$, for all $p, p' \in (1/2, 1)$.

Moreover, $\vdash_{(p} \; \subseteq \; \vdash_{[p}$ and $\vdash^f_{[p} \; = \; \vdash^f_{(p}$ for all $p \in (1/2, 1)$.


\item[P2.] $\vdash_{[n} \; = \; \vdash_{[n'}$ and $\vdash_{(n} \; = \; \vdash_{(n'}$, for all $n, n' \in (0, 1/2)$.

Moreover, $\vdash_{(n} \; \subseteq \; \vdash_{[n}$ and $\vdash^f_{[n} \; = \; \vdash^f_{(n}$ for all $n \in (0, 1/2)$.

\item[P3.] $\vdash_{[p} ~\varsubsetneq ~\vdash_1$,  for any $p \in (1/2,1)$.

\item[P4.] $\vdash_1$ and ~$\vdash_{[1/2}$ are not comparable. 
\item[P5.] $\vdash_{[p}$ and ~$\vdash_{[1/2}$, as well as $\vdash^f_{[p}$ and ~$\vdash_{[1/2}$, are not comparable, for any $p \in (1/2,1)$. 
\item[P6.] $\vdash^f_{[p}  ~\varsubsetneq  ~\vdash_{(1/2}$,  for any $p \in (1/2,1)$.
\item[P7.] $\vdash_{[p}$ and ~$\vdash_{[n}$ are not comparable, for any $p \in (1/2,1)$ and any $n \in (0,1/2)$. The same holds for $\vdash^f_{[p}$ and ~$\vdash^f_{[n}$. 
\item[P8.]  $\vdash_{[n} ~\varsubsetneq ~\vdash_{[1/2}$,  for any $n \in (0,1/2)$. 
\item[P9.] $\vdash_{(0}$, $\vdash_{[1/2}$ and $\vdash_{(1/2}$ are not pair-wise comparable. 
\item[P10.]  $\vdash^f_{[n}  \; \varsubsetneq  \; \vdash_{(0}$, for any $n \in (0, 1/2)$.

\end{enumerate}
\end{proposition}

\begin{proof}  \mbox{}

\begin{enumerate}
\item[P1.]  We divide the proof in four steps:

(i) That $\vdash_{[p} \; = \; \vdash_{[p'}$ and $\vdash_{(p} \; = \; \vdash_{(p'}$ is an easy consequence of the fact that for every $p, p' \in (1/2,1)$  it is possible to define an automorphism $f$ of $[0,1]_{G\sim}$ such that $f(p) = p'$. Let us then show that  $\vdash^f_{[p} \; =\;  \vdash^f_{(p}$ for every $p \in (1/2, 1)$.

(ii) Assume $\{\varphi_i : i \in I\} \vdash^f_{[p} \psi$, with $I$ finite, for some $p \in (1/2, 1)$. Let $q$ such that $1/2 < q < p$, and let $e$ be an evaluation such that $e(\varphi_i) > q $ for all $i \in I$. 
Let $p' = \min_{i \in I} e(\varphi_i)$. Obviously $p' > q$. Then, by (i), we also have $\{\varphi_i : i \in I\} \vdash^f_{[p'} \psi$, and therefore we have $e(\psi) \geq p' > q$, and hence $\{\varphi_i :  i \in I\} \vdash^f_{(q} \psi$. Therefore, we have $\vdash^f_{[p}  \; \subseteq  \; \vdash^f_{(q}$ for all $1/2 < q < p$.

(iii) Recall from (i) that $\Gamma \vdash_{(p} \varphi$ iff $\Gamma \vdash_{(p'} \varphi$ for all $1/2 < p' < 1$. Let $p_1, p_2, \ldots, p_n, \ldots $ be an increasing sequence of values $p_i \in (1/2, p)$ such that $\lim_n p_n = p$. 
%
Suppose $\Gamma \vdash_{(p} \varphi$, and further assume $e(\psi) \geq p$ for all $\psi \in \Gamma$. Clearly, for each $p_i$, $e(\psi) > p_i$ for all $\psi \in \Gamma$. Since $\Gamma \vdash_{(p_i} \varphi$ for each $p_i$, we have that $e(\varphi) > p_i$ for each $p_i$.  Hence $e(\varphi) \geq p$. 

(iv) Assume $\{\varphi_i :  i \in I\} \vdash^f_{(q} \psi$, with $I$ finite, for some $q \in (1/2, 1)$. Let $p$ be such that $ q < p < 1$, and let $e$ be an evaluation such that $e(\varphi_i) \geq p $ for all $i \in I$. 
Let $q' = \min_{i \in I} e(\varphi_i)$. Obviously $q' \geq p$. Then, by (i), we also have $\{\varphi_i :  i \in I\} \vdash^f_{(q'} \psi$, and therefore we have $e(\psi) \geq q' \geq p$, and hence $\{\varphi_i :  i \in I\} \vdash^f_{p} \psi$. Therefore, we have $\vdash^f_{(q}  \; \subseteq  \; \vdash^f_{[p}$ for all $1/2 < q < p$.

\item[P2.] The proofs are analogous to those of P1.

\item[P3.] Assume $\{\varphi_i :  i \in I\} \vdash_{[p} \psi$ for a given $p  \in (1/2,1)$, and let $e$ be an evaluation such that $e(\varphi_i) = 1$ for all $i \in I$. Since it is also true that $e(\varphi_i) \geq p'$ for all $p'  \in (1/2,1)$, by P1 it follows that $\{\varphi_i :  i \in I\} \vdash_{[p'} \psi$ for all $p'  \in (1/2,1)$, and hence $e(\psi) \geq p'$ for all $p'  \in (1/2,1)$, and thus $e(\psi) = 1$. Therefore $\{\varphi_i :  i \in I\} \vdash_1 \psi$.

The strict inclusion can be easily noticed since, e.g. it holds that  $\varphi \vdash_1 \Delta \varphi$ but $\varphi \nvdash_{[p} \Delta \varphi$ for any $p < 1$.


\item[P4.] It clearly holds that, on the one hand, $\Delta(\varphi \leftrightarrow \ninv \varphi) \vdash_{[1/2} \varphi$ but $\Delta(\varphi \leftrightarrow \ninv \varphi) \nvdash_{1} \varphi$, while on the other hand,
$\varphi \vdash_1 \Delta \varphi$ but $\varphi \nvdash_{[1/2} \Delta \varphi$

\item[P5.]  It follows from noticing that $\Delta(\varphi \leftrightarrow \ninv \varphi) \land \varphi \vdash_{[p} \bot$ and $\Delta(\varphi \leftrightarrow \ninv \varphi) \land \varphi \nvdash_{[1/2} \bot$, while
$\Delta(\varphi \leftrightarrow \ninv \varphi) \vdash_{[1/2} \varphi$ and $\Delta(\varphi \leftrightarrow \ninv \varphi) \nvdash_{[p} \varphi$.

\item[P6.]
Assume that, for a given $p  \in (1/2,1)$, $\{\varphi_i :  i \in I\} \vdash_{[p} \psi$, with $I$ finite, and let $e$ be an evaluation such that $e(\varphi_i) > 1/2 $ for all $i \in I$. 
Let $p' = \min_{i \in I} e(\varphi_i)$. Obviously $p' > 1/2$. Then, from P1 we also have $\{\varphi_i :  i \in I\} \vdash_{[p'} \psi$, and therefore we have $e(\psi) \geq p' > 1/2$, and hence $\{\varphi_i :  i \in I\} \vdash_{(1/2} \psi$. Therefore, we have $\vdash_{[p}  \; \subseteq  \; \vdash_{(1/2}$.

That the inclusion is strict follows from observing that $\ninv \Delta (\varphi \to \ninv \varphi) \vdash_{(1/2} \varphi$, but $ \ninv \Delta (\varphi \to \ninv \varphi) \nvdash_{[p} \varphi$.

\item[P7.]  It follows from observing (i) $\Delta(\varphi \leftrightarrow \ninv \varphi) \vdash_{[n} \varphi$ and $\Delta(\varphi \leftrightarrow \ninv \varphi) \nvdash_{[p} \varphi$, and  (ii)
$ \varphi \vdash_{[p} \ninv \Delta (\varphi \to  \ninv \varphi)$ and $\varphi \nvdash_{[n} \ninv \Delta (\varphi \to \ninv \varphi)$.

\item[P8.]  That $\vdash_{[n} ~\subseteq ~\vdash_{[1/2}$  is proved in a similar way to P3. The strict inclusion is a consequence of the following facts:

(i) $\varphi \land \ninv \varphi \vdash_{[1/2} \Delta(\varphi \leftrightarrow \ninv \varphi)$

(ii) $\varphi \land \ninv \varphi \not\vdash_{[n} \Delta(\varphi \leftrightarrow \ninv \varphi)$

Notice that $e(\varphi \land \ninv \varphi) \geq 1/2$ iff $e(\varphi) = 1/2$ iff $e(\varphi \leftrightarrow \ninv \varphi) = 1$, while if $e(\varphi) = n$, then $e(\varphi \land \ninv \varphi) = e(\varphi \leftrightarrow \ninv \varphi) = n$, but $e(\Delta(\varphi \leftrightarrow \ninv \varphi)) = 0$.   



\item[P9.]  That $\vdash_{[1/2}$ and $\vdash_{(1/2}$ are not comparable results from noticing e.g. (i) $\Delta(\varphi \leftrightarrow \ninv \varphi) \vdash_{[1/2} \varphi$ but $\Delta(\varphi \leftrightarrow \ninv \varphi) \nvdash_{(1/2} \varphi$, and (ii)
$\varphi \vdash_{(1/2} \ninv \Delta(\varphi \to \ninv \varphi)$ but $\varphi \nvdash_{[1/2} \ninv \Delta(\varphi \to \ninv \varphi)$.

On the other hand, it is easy to check that $\bot$ follows from $\varphi \land \Delta(\varphi \to \ninv \varphi) \land \neg \Delta (\ninv \varphi \to \varphi)$ in $\vdash_{(1/2}$ and $\vdash_{[1/2}$, but not in $\vdash_{(0}$. Conversely, 
$\neg\neg \varphi \land \neg \Delta \varphi  
\vdash_{(0} \varphi \land \ninv \varphi$, but this is neither the case for $\vdash_{(1/2}$ nor for $\vdash_{[1/2}$. 


\item[P10.]  Assume $\{\varphi_i :  i \in I\} \vdash_{[n} \psi$ for a given $n  \in (0, 1/2)$ and a finite set $I$, and let $e$ be an evaluation such that $e(\varphi_i) > 0 $ for all $i \in I$. 
Let $n' = \min_{i \in I} e(\varphi_i)$. Obviously $n' > 0$ and $e(\varphi_i) \geq n'$, for all $i \in I$. Then, from P1, we also have that $\{\varphi_i :  i \in I\} \vdash_{[n'} \psi$, and hence we have $e(\psi) \geq n' > 0$.  This means $\{\varphi_i :  i \in I\} \vdash_{(0} \psi$. Therefore, we have $\vdash^f_{[n}  \; \subseteq  \; \vdash_{(0}$.

On the other hand, $\neg \neg\varphi \vdash_{(0}  \varphi$ but  $\neg \neg\varphi \nvdash_{[n}  \varphi$, hence we have proved that $\vdash^f_{[n}  \; \varsubsetneq  \; \vdash_{(0}$.
\end{enumerate}
\end{proof}

\end{document}
